% Glosario / nomenclatura. La plantilla UMA proporciona el comando
% \nomenclature{}{} y \printnomenclature. La fuente única de los
% términos vive en docs/glossary.md; aquí los transcribimos al formato
% LaTeX. Cuando la terminología cambie en glossary.md, hay que
% propagarlo aquí (apuntado como deuda en claude/ideas-pendientes.md).

% --- Árbol de descriptores ---
\nomenclature{\textbf{Descriptor}}{Objeto del orquestador que representa
  una VM (definición + estado). Analogía con un descriptor de fichero en
  un sistema operativo.}
\nomenclature{\textbf{Carpeta}}{Nodo agrupador del árbol, contenedor de
  descriptores y de otras carpetas. Soporta importación explícita de
  recursos heredables.}
\nomenclature{\textbf{Path materializado}}{Identificador único de una
  carpeta o descriptor; cadena con la ruta completa desde la raíz
  (e.g.\ \texttt{/lab/networks/student01}).}

% --- Recursos heredables ---
\nomenclature{\textbf{Plantilla (template)}}{Recurso declarable con
  \texttt{define templates:}; ofrece valores por defecto reutilizables
  por descriptores que la referencian con \texttt{use template:}.}
\nomenclature{\textbf{Hipervisor (recurso)}}{Recurso declarable con
  \texttt{define hypervisors:}; nombra una instancia de hipervisor
  concreta (con su tipo y datos de acceso) referenciable por
  descriptores y plantillas con \texttt{use hypervisor:}.}
\nomenclature{\textbf{Importación}}{Mecanismo por el que una carpeta
  declara qué recursos del padre quiere hacer visibles aquí; sin
  importación, el recurso no es accesible desde esta rama.}
\nomenclature{\textbf{Cierre léxico}}{Regla por la que las referencias
  internas de una plantilla se resuelven desde donde la plantilla está
  definida, no desde donde se usa.}

% --- Estado ---
\nomenclature{\textbf{State}}{Eje de estado del descriptor desde el
  punto de vista del orquestador: \texttt{provisioned}, \texttt{deployed}
  (con o sin \texttt{drifted}), \texttt{broken}, \texttt{unreachable}.}
\nomenclature{\textbf{Runtime}}{Eje de estado de la VM desde el punto de
  vista del hipervisor: \texttt{running}, \texttt{stopped},
  \texttt{paused}, \texttt{unknown}.}

% --- Operaciones ---
\nomenclature{\textbf{Job}}{Registro persistido de una operación, con su
  identificador único, tipo, objetivo, estado y resultado. Garantiza
  auditoría y prepara la asincronía futura.}
\nomenclature{\textbf{Command}}{Patrón de diseño aplicado a las
  operaciones: cada acción encapsula \texttt{validate()} y
  \texttt{execute()}, permitiendo \emph{dry-run} (\texttt{plan}),
  ejecución diferida y reintento.}
\nomenclature{\textbf{Conector}}{Adaptador entre el orquestador y un
  hipervisor concreto. Implementa el contrato ABC
  \texttt{HypervisorConnector}.}

% --- Acciones del usuario ---
\nomenclature{\textbf{load}}{Acción que carga un documento YAML de
  definiciones en una carpeta destino.}
\nomenclature{\textbf{deploy / undeploy}}{Crear o destruir la VM en el
  hipervisor a partir del descriptor (transición \texttt{provisioned} a
  \texttt{deployed} y viceversa).}
\nomenclature{\textbf{start / stop}}{Encender o apagar la VM ya
  desplegada (afecta al runtime, no al state).}

\printnomenclature
\clearpage
